% =========================================================
% Structure of the Rational Numbers
% =========================================================

\subsection{Structure of the Rational Numbers}

\begin{remark}[Structural overview]

The rational numbers $\mathbb{Q}$ possess a surprisingly rich internal
structure.  Several seemingly different constructions reveal the same
underlying organization of rational numbers and their approximating
power.

\end{remark}

% ---------------------------------------------------------
% Construction
% ---------------------------------------------------------

\begin{tcolorbox}[colback=gray!6, colframe=gray!40, arc=2pt,
title={Construction of $\mathbb{Q}$}]

\[
\mathbb{Q}
=
\{(a,b)\in\mathbb{Z}\times\mathbb{Z}\;:\; b\neq0\}/\sim
\]

where

\[
(a,b)\sim(c,d)
\quad\Longleftrightarrow\quad
ad=bc.
\]

Each rational number is an equivalence class

\[
[(a,b)].
\]

\end{tcolorbox}

\medskip

% ---------------------------------------------------------
% Field structure
% ---------------------------------------------------------

\begin{tcolorbox}[colback=gray!6, colframe=gray!40, arc=2pt,
title={Algebraic Structure}]

\begin{itemize}

\item $(\mathbb{Q},+,\cdot)$ forms a \textbf{field}.

\item $\mathbb{Q}$ is an \textbf{ordered field} with order inherited
from $\mathbb{Z}$.

\item The \textbf{Archimedean property} holds:

\[
\forall x\in\mathbb{Q}
\quad
\exists n\in\mathbb{N}
\quad
n>x.
\]

\end{itemize}

\end{tcolorbox}

\medskip

% ---------------------------------------------------------
% Density
% ---------------------------------------------------------

\begin{tcolorbox}[colback=gray!6, colframe=gray!40, arc=2pt,
title={Density of $\mathbb{Q}$}]

Between any two rational numbers there exists another rational number.

\[
\forall a,b\in\mathbb{Q},\ a<b
\quad
\exists r\in\mathbb{Q}
\quad
a<r<b.
\]

A simple construction is the \textbf{mediant}

\[
\frac{a}{b},\frac{c}{d}
\quad\Rightarrow\quad
\frac{a+c}{b+d}.
\]

\end{tcolorbox}

\medskip

% ---------------------------------------------------------
% Farey sequences
% ---------------------------------------------------------

\begin{tcolorbox}[colback=gray!6, colframe=gray!40, arc=2pt,
title={Farey Sequences}]

The Farey sequence $F_n$ contains all reduced fractions

\[
\frac{a}{b}
\]

with

\[
0\le a\le b\le n
\]

listed in increasing order.

If

\[
\frac{a}{b},\frac{c}{d}
\]

are neighbors in a Farey sequence, then

\[
bc-ad=1.
\]

Their spacing satisfies

\[
\frac{c}{d}-\frac{a}{b}=\frac{1}{bd}.
\]

\end{tcolorbox}

\medskip

% ---------------------------------------------------------
% Stern–Brocot tree
% ---------------------------------------------------------

\begin{tcolorbox}[colback=gray!6, colframe=gray!40, arc=2pt,
title={Stern--Brocot Tree}]

All positive rational numbers appear exactly once in the
Stern–Brocot tree.

Starting with

\[
\frac{0}{1},\quad\frac{1}{0},
\]

new fractions are generated by mediants:

\[
\frac{a}{b},\frac{c}{d}
\quad\mapsto\quad
\frac{a+c}{b+d}.
\]

This process recursively enumerates every positive rational number.

\end{tcolorbox}

\medskip

% ---------------------------------------------------------
% Rational approximation
% ---------------------------------------------------------

\begin{tcolorbox}[colback=gray!6, colframe=gray!40, arc=2pt,
title={Rational Approximation}]

Rational numbers approximate real numbers extremely well.

\textbf{Dirichlet's Approximation Theorem}

For every real number $x$ there exist infinitely many rational numbers

\[
\frac{p}{q}
\]

such that

\[
\left|x-\frac{p}{q}\right|
<
\frac{1}{q^{2}}.
\]

\end{tcolorbox}

\medskip

% ---------------------------------------------------------
% Gaps
% ---------------------------------------------------------

\begin{tcolorbox}[colback=gray!6, colframe=gray!40, arc=2pt,
title={Gaps in $\mathbb{Q}$}]

Despite their density, the rational numbers are not complete.

For example,

\[
x^{2}=2
\]

has no rational solution.

Thus there exist sequences of rational numbers that approach
limits which are not rational.

\end{tcolorbox}

\medskip

% ---------------------------------------------------------
% Completion
% ---------------------------------------------------------

\begin{tcolorbox}[colback=gray!6, colframe=gray!40, arc=2pt,
title={Completion to the Real Numbers}]

The real numbers can be constructed as equivalence classes
of Cauchy sequences of rational numbers:

\[
\mathbb{R}
=
\{\text{Cauchy sequences in }\mathbb{Q}\}/\sim.
\]

Two sequences are equivalent if their difference converges to $0$.

Thus the real numbers complete the rational numbers by filling
the gaps.

\end{tcolorbox}

\medskip

\begin{remark}[Conceptual picture]

\[
\mathbb{Z}
\quad\longrightarrow\quad
\mathbb{Q}
\quad\longrightarrow\quad
\mathbb{R}
\]

\[
\text{integers}
\quad
\text{fractions}
\quad
\text{limits}
\]

The rational numbers form a dense ordered field with a rich
combinatorial structure, but they are not complete.  The real numbers
arise by adjoining the limits that are missing from $\mathbb{Q}$.

\end{remark}